\section{Measured structure: current, rotor reference and speed}
\label{sec:experimental-results}
\subsection{Support counts and the distinction between a pair and a parameter}
\Cref{tab:experimental-summary} reports both stages of current aggregation.
The outlier slices contain 12882, 12957 and 12959 raw samples, respectively;
each produces 4220 operating-index means and 4154 requested keys. Among the
keys, 48 have zero q current, 2344 form 1172 pairs, and 1762 nonzero-q keys
remain unmatched. The relative minimum-current rule leaves 1002 equivalent
residuals. These support losses are not filled by interpolation.
Each multi-speed slice has 327 samples, 109 operating-index means, 104 keys,
eight zero-q keys, 39 pairs and 18 unmatched nonzero-q keys. Three pairs per
slice are excluded only from division by $a_q$, leaving 36 valid pairs.

\begin{table}[tb]
\centering\small
\begin{tabular}{@{}rrrrrrrrr@{}}
\toprule
$n$ & $T_{\rm rot,ref}$ & $N_{\rm op}$ & $N_{\rm key}$ & $N_{\rm pair}$ & $N_{\rm valid}$ & $N_{\rm high}$ & Median & MAD\\
rpm & $^\circ$C & & & & & & m$\Omega$ & m$\Omega$\\
\midrule
2000 & 40 & 4220 & 4154 & 1172 & 1002 & 79 & 0.096 & 0.128 \\
2000 & 60 & 4220 & 4154 & 1172 & 1002 & 80 & -0.579 & 0.080 \\
2000 & 80 & 4220 & 4154 & 1172 & 1002 & 79 & -1.268 & 0.160 \\
\midrule
1000 & 70 & 109 & 104 & 39 & 36 & 6 & -2.157 & 0.189 \\
3000 & 70 & 109 & 104 & 39 & 36 & 7 & -5.334 & 1.155 \\
\bottomrule
\end{tabular}

\caption{Reproduced slice counts and high-current $\Delta R_{\mathrm{eq}}$
summaries. $N_{\rm op}$ counts operating-index groups; $N_{\rm key}$ counts
aggregated requested points. Valid means $a_q/a_{q,\max}>0.10$; high means
$a_q/a_{q,\max}\geq0.85$. Medians and raw MAD describe the high subset only.}
\label{tab:experimental-summary}
\end{table}

\subsection{Constant-speed temperature-reference experiment}
\Cref{fig:experimental-flux} displays both reconstructed components at all
measured operating-index means. Their broad nonlinear shapes alone do not
identify the origin of a small asymmetry. Pairing exposes the complementary
structure in \cref{fig:experimental-residuals}. Both forbidden components
are appreciable and spatially structured; retaining $r_q$ avoids reducing
the diagnosis to a scalar resistance interpretation.
Across all 1172 pairs the RMS values $(r_d,r_q)$ are
$(2.013,1.191)$, $(1.746,1.309)$ and $(1.728,1.690)$ mWb at 40, 60 and
80 $^\circ$C. These RMS values include the small-q pairs excluded from
the equivalent transformation, and describe map variation rather than
calibrated measurement noise.
% Generated by python/export_experiments.py; do not edit manually.
\begin{figure}[p]
\centering
\begin{tikzpicture}
\begin{groupplot}[group style={group size=2 by 3,horizontal sep=32mm,vertical sep=19mm},width=.40\textwidth,height=4.8cm]
\nextgroupplot[sample map axis,width=.35\textwidth,height=4.8cm,title={$40$ $^\circ$C; $\psi_d$ / mWb},xlabel={$i_d$ / A},ylabel={$i_q$ / A},point meta min=-108.825209645,point meta max=108.825209645]
\addplot[sample map] table[x=id_plot,y=iq_plot,meta=psi_d_mWb]{figures/data/experimental/outlier_2000_40_flux.dat};
\nextgroupplot[sample map axis,width=.35\textwidth,height=4.8cm,title={$40$ $^\circ$C; $\psi_q$ / mWb},xlabel={$i_d$ / A},ylabel={$i_q$ / A},point meta min=-225.768924961,point meta max=225.768924961]
\addplot[sample map] table[x=id_plot,y=iq_plot,meta=psi_q_mWb]{figures/data/experimental/outlier_2000_40_flux.dat};
\nextgroupplot[sample map axis,width=.35\textwidth,height=4.8cm,title={$60$ $^\circ$C; $\psi_d$ / mWb},xlabel={$i_d$ / A},ylabel={$i_q$ / A},point meta min=-108.825209645,point meta max=108.825209645]
\addplot[sample map] table[x=id_plot,y=iq_plot,meta=psi_d_mWb]{figures/data/experimental/outlier_2000_60_flux.dat};
\nextgroupplot[sample map axis,width=.35\textwidth,height=4.8cm,title={$60$ $^\circ$C; $\psi_q$ / mWb},xlabel={$i_d$ / A},ylabel={$i_q$ / A},point meta min=-225.768924961,point meta max=225.768924961]
\addplot[sample map] table[x=id_plot,y=iq_plot,meta=psi_q_mWb]{figures/data/experimental/outlier_2000_60_flux.dat};
\nextgroupplot[sample map axis,width=.35\textwidth,height=4.8cm,title={$80$ $^\circ$C; $\psi_d$ / mWb},xlabel={$i_d$ / A},ylabel={$i_q$ / A},point meta min=-108.825209645,point meta max=108.825209645]
\addplot[sample map] table[x=id_plot,y=iq_plot,meta=psi_d_mWb]{figures/data/experimental/outlier_2000_80_flux.dat};
\nextgroupplot[sample map axis,width=.35\textwidth,height=4.8cm,title={$80$ $^\circ$C; $\psi_q$ / mWb},xlabel={$i_d$ / A},ylabel={$i_q$ / A},point meta min=-225.768924961,point meta max=225.768924961]
\addplot[sample map] table[x=id_plot,y=iq_plot,meta=psi_q_mWb]{figures/data/experimental/outlier_2000_80_flux.dat};
\end{groupplot}
\end{tikzpicture}
\caption{Reconstructed measured flux samples at 2000 rpm. Columns show $\psi_d$ and $\psi_q$; rows show rotor references 40, 60 and 80 $^\circ$C. Color denotes flux in mWb with one fixed range per component across temperatures. Every operating-point mean is shown; no surface fit or interpolation is applied.}
\label{fig:experimental-flux}
\end{figure}

% Generated by python/export_experiments.py; do not edit manually.
\begin{figure}[p]
\centering
\begin{tikzpicture}
\begin{groupplot}[group style={group size=2 by 3,horizontal sep=32mm,vertical sep=19mm},width=.40\textwidth,height=4.8cm]
\nextgroupplot[sample map axis,width=.35\textwidth,height=4.8cm,title={$40$ $^\circ$C; $r_d$ / mWb},xlabel={$i_{d,\mathrm{key}}$ / A},ylabel={$a_q$ / A},point meta min=-4.39315188828,point meta max=4.39315188828]
\addplot[sample map] table[x=id_plot,y=iq_plot,meta=rd_mWb]{figures/data/experimental/outlier_2000_40_pairs.dat};
\nextgroupplot[sample map axis,width=.35\textwidth,height=4.8cm,title={$40$ $^\circ$C; $r_q$ / mWb},xlabel={$i_{d,\mathrm{key}}$ / A},ylabel={$a_q$ / A},point meta min=-6.36492032923,point meta max=6.36492032923]
\addplot[sample map] table[x=id_plot,y=iq_plot,meta=rq_mWb]{figures/data/experimental/outlier_2000_40_pairs.dat};
\nextgroupplot[sample map axis,width=.35\textwidth,height=4.8cm,title={$60$ $^\circ$C; $r_d$ / mWb},xlabel={$i_{d,\mathrm{key}}$ / A},ylabel={$a_q$ / A},point meta min=-4.39315188828,point meta max=4.39315188828]
\addplot[sample map] table[x=id_plot,y=iq_plot,meta=rd_mWb]{figures/data/experimental/outlier_2000_60_pairs.dat};
\nextgroupplot[sample map axis,width=.35\textwidth,height=4.8cm,title={$60$ $^\circ$C; $r_q$ / mWb},xlabel={$i_{d,\mathrm{key}}$ / A},ylabel={$a_q$ / A},point meta min=-6.36492032923,point meta max=6.36492032923]
\addplot[sample map] table[x=id_plot,y=iq_plot,meta=rq_mWb]{figures/data/experimental/outlier_2000_60_pairs.dat};
\nextgroupplot[sample map axis,width=.35\textwidth,height=4.8cm,title={$80$ $^\circ$C; $r_d$ / mWb},xlabel={$i_{d,\mathrm{key}}$ / A},ylabel={$a_q$ / A},point meta min=-4.39315188828,point meta max=4.39315188828]
\addplot[sample map] table[x=id_plot,y=iq_plot,meta=rd_mWb]{figures/data/experimental/outlier_2000_80_pairs.dat};
\nextgroupplot[sample map axis,width=.35\textwidth,height=4.8cm,title={$80$ $^\circ$C; $r_q$ / mWb},xlabel={$i_{d,\mathrm{key}}$ / A},ylabel={$a_q$ / A},point meta min=-6.36492032923,point meta max=6.36492032923]
\addplot[sample map] table[x=id_plot,y=iq_plot,meta=rq_mWb]{figures/data/experimental/outlier_2000_80_pairs.dat};
\end{groupplot}
\end{tikzpicture}
\caption{Forbidden components at all 1172 matched pairs per rotor-reference slice. Color scales are fixed per component across rows and include the full extrema. The horizontal label is the requested d-current key; $a_q$ is the mean measured q-current magnitude of the pair. Local extrema and unmeasured gaps are retained.}
\label{fig:experimental-residuals}
\end{figure}


The constant-resistance hypothesis becomes easier to assess after removing
the current/speed factors. Over the 1002 valid pairs,
the equivalent residual ranges from $-1.595$ to $13.989$ m$\Omega$ at
40 $^\circ$C, $-3.370$ to $12.667$ m$\Omega$ at 60 $^\circ$C, and
$-3.389$ to $11.965$ m$\Omega$ at 80 $^\circ$C.
The respective full-region medians are $1.680$, $0.967$ and
$0.250$ m$\Omega$, with raw MAD $1.576$, $1.531$ and $1.553$ m$\Omega$.
\Cref{fig:experimental-equivalent,fig:high-current} retain this nonconstant
structure and all valid extrema. A constant shift cannot explain it globally.
% Generated by python/export_experiments.py; do not edit manually.
\begin{figure}[p]
\centering
\begin{tikzpicture}
\begin{groupplot}[group style={group size=2 by 2,horizontal sep=32mm,vertical sep=22mm},width=.40\textwidth,height=5.2cm]
\nextgroupplot[sample map axis,width=.35\textwidth,height=4.8cm,title={$40$ $^\circ$C; $\Delta R_{\mathrm{eq}}$ / m$\Omega$},xlabel={$i_{d,\mathrm{key}}$ / A},ylabel={$a_q$ / A},point meta min=-13.9889154009,point meta max=13.9889154009]
\addplot[sample map] table[x=id_plot,y=iq_plot,meta=req_mOhm]{figures/data/experimental/outlier_2000_40_valid.dat};
\nextgroupplot[sample map axis,width=.35\textwidth,height=4.8cm,title={$60$ $^\circ$C; $\Delta R_{\mathrm{eq}}$ / m$\Omega$},xlabel={$i_{d,\mathrm{key}}$ / A},ylabel={$a_q$ / A},point meta min=-13.9889154009,point meta max=13.9889154009]
\addplot[sample map] table[x=id_plot,y=iq_plot,meta=req_mOhm]{figures/data/experimental/outlier_2000_60_valid.dat};
\nextgroupplot[sample map axis,width=.35\textwidth,height=4.8cm,title={$80$ $^\circ$C; $\Delta R_{\mathrm{eq}}$ / m$\Omega$},xlabel={$i_{d,\mathrm{key}}$ / A},ylabel={$a_q$ / A},point meta min=-13.9889154009,point meta max=13.9889154009]
\addplot[sample map] table[x=id_plot,y=iq_plot,meta=req_mOhm]{figures/data/experimental/outlier_2000_80_valid.dat};
\nextgroupplot[diagnostic axis,width=.35\textwidth,height=4.8cm,title={High-current summary},xlabel={Rotor reference / $^\circ$C},ylabel={Median / m$\Omega$},xmin=30,xmax=90]
\addplot[reference samples,error bars/.cd,y dir=both,y explicit] coordinates {(40.0,0.09585958213328295) +- (0,0.1278655654395009) (60.0,-0.5790103359149268) +- (0,0.08035901637329465) (80.0,-1.2684063668812) +- (0,0.15955533200440164)};
\end{groupplot}
\end{tikzpicture}
\caption{Is the resistance-equivalent residual spatially constant? All valid pairs are shown at a common, unclipped color range in m$\Omega$. The last panel summarizes $a_q/a_{q,\max}\geq0.85$ with median and raw MAD bars, which represent descriptive spread rather than confidence intervals. Rotor reference is not winding temperature.}
\label{fig:experimental-equivalent}
\end{figure}

% Generated by python/export_experiments.py; do not edit manually.
\begin{figure}[p]
\centering
\begin{tikzpicture}
\begin{axis}[diagnostic axis,width=.85\textwidth,height=7cm,xlabel={$a_q/a_{q,\max}$},ylabel={$\Delta R_{\mathrm{eq}}$ / m$\Omega$},title={Retain the spatial spread},legend columns=3,xmin=.1,xmax=1.02]
\addplot[reference samples] table[x=relative_iq,y=req_mOhm]{figures/data/experimental/outlier_2000_40_valid.dat};
\addlegendentry{40 $^\circ$C}
\addplot[measured samples] table[x=relative_iq,y=req_mOhm]{figures/data/experimental/outlier_2000_60_valid.dat};
\addlegendentry{60 $^\circ$C}
\addplot[only marks,derived quantity,mark=triangle*] table[x=relative_iq,y=req_mOhm]{figures/data/experimental/outlier_2000_80_valid.dat};
\addlegendentry{80 $^\circ$C}
\draw[guide quantity,guide pattern] (rel axis cs:.8152173913,0) -- (rel axis cs:.8152173913,1);
\end{axis}
\end{tikzpicture}
\caption{Each marker is one valid matched pair, projected onto relative measured q-current. The narrower spread near maximum q-current is an empirical feature of this dataset. The guide marks the initial 0.85 summary threshold; it does not define a universal identification region. No outlier removal is used.}
\label{fig:high-current}
\end{figure}


The high-current subset is flatter: its median/raw-MAD values are
$0.096/0.128$, $-0.579/0.080$ and $-1.268/0.160$ m$\Omega$, with
79, 80 and 79 pairs. The high-current ranges are only
$[-0.174,0.461]$, $[-0.851,-0.239]$ and $[-1.661,-0.880]$ m$\Omega$.
Thus a local constant equivalent is a useful descriptive compression of
part of this experiment, while the rest of the map remains unexplained.
The supplied code also reports $\hat R_s-\operatorname{median}
(\Delta R_{\mathrm{eq}})$: $7.599$, $8.274$ and $8.963$ m$\Omega$.
These are conditional resistance equivalents, not identified true winding
resistances or recommended corrections.

The shift with rotor reference is observed, but its thermal cause is not
isolated. \texttt{Rotor\_temp\_ref} is not winding temperature. An additional
recorded stator-temperature \emph{reference},
\texttt{PE\_Char\_\allowbreak Temp\_Stator\_\allowbreak ref\_DegCel}, happens to equal 40, 60 and
80 $^\circ$C in the respective blocks; it is unused in the current method
and likewise does not verify the conductor temperature. Actual stator thermal
signals, their sensor meaning, settling and acquisition history should be
checked before attributing the shift quantitatively to $R_s(T)$.
The recorded \texttt{PE\_speed\_rpm} medians are approximately 1999.14,
1999.23 and 1999.32 rpm, but individual samples vary. The reconstruction
deliberately retains the requested 2000-rpm slice value.

\subsection{Constant rotor-reference speed experiment}
At a 70-$^\circ$C rotor reference, high-current medians differ from
$-2.157$ m$\Omega$ at 1000 rpm to $-5.334$ m$\Omega$ at 3000 rpm.
Raw MAD increases from $0.189$ to $1.155$ m$\Omega$; only six and seven
pairs contribute. Corresponding conditional resistance equivalents are
$12.157$ and $15.334$ m$\Omega$ relative to the common used
$10$ m$\Omega$. They cannot both identify one constant true resistance
under $H_0$.
The full valid region is also nonconstant: ranges are
$[-2.431,33.912]$ and $[-7.541,28.267]$ m$\Omega$;
medians/raw-MAD are $0.269/2.108$ and $-1.444/1.568$ m$\Omega$.
RMS $(r_d,r_q)$ over all 39 pairs is $(0.385,0.974)$ mWb at 1000 rpm
and $(0.206,0.899)$ mWb at 3000 rpm. Sparse matched support and differing
high-current membership limit the interpretation of slice medians.
% Generated by python/export_experiments.py; do not edit manually.
\begin{figure}[p]
\centering
\begin{tikzpicture}
\begin{groupplot}[group style={group size=2 by 1,horizontal sep=18mm},diagnostic axis,height=6cm]
\nextgroupplot[title={Speed and operating-point dependence},xlabel={$a_q/a_{q,\max}$},ylabel={$\Delta R_{\mathrm{eq}}$ / m$\Omega$},legend columns=2]
\addplot[reference samples] table[x=relative_iq,y=req_mOhm]{figures/data/experimental/multi_1000_70_valid.dat};
\addlegendentry{1000 rpm}
\addplot[measured samples] table[x=relative_iq,y=req_mOhm]{figures/data/experimental/multi_3000_70_valid.dat};
\addlegendentry{3000 rpm}
\nextgroupplot[sample map axis,width=.38\textwidth,height=6cm,title={Speed difference / m$\Omega$},xlabel={$i_{d,\mathrm{key}}$ / A},ylabel={$|i_{q,\mathrm{key}}|$ / A},point meta min=-15.178282324759852,point meta max=15.178282324759852]
\addplot[sparse sample map] table[x=id_key,y=iq_key,meta=difference_mOhm]{figures/data/experimental/common_speed.dat};
\end{groupplot}
\end{tikzpicture}
\caption{All 36 valid equivalent-residual pairs at each speed (left), and differences $\Delta R_{\mathrm{eq},3000}-\Delta R_{\mathrm{eq},1000}$ on common keys (right). Separate high-current summaries contain only six and seven pairs. Requested-key matching preserves nominal geometry but does not force identical measured currents.}
\label{fig:experimental-speed-structure}
\end{figure}


\subsection{A direct test at common requested points}
The valid-key inner join contains 36 common pairs. This comparison avoids
explaining speed differences by comparing different requested geometries.
Under a sole constant resistance mismatch and equal measured pair currents,
\begin{equation}
 r_{d,3000}^{(H_0)}=\frac{1000}{3000}r_{d,1000}
 =\frac13r_{d,1000}.\label{eq:direct-speed-test}
\end{equation}
\Cref{fig:experimental-scaling} compares that prediction directly with the
observed residual. The RMS of $r_{d,3000}-r_{d,1000}/3$ is
\SI{0.1573}{\milli\weber}, and its median is
\SI{0.1040}{\milli\weber}. For context, the observed high-speed residual
RMS on these same 36 points is \SI{0.2135}{\milli\weber}.
The discrepancy is therefore not hidden by the speed transformation.
% Generated by python/export_experiments.py; do not edit manually.
\begin{figure}[p]
\centering
\begin{tikzpicture}
\begin{groupplot}[group style={group size=2 by 1,horizontal sep=18mm},diagnostic axis,height=6cm]
\nextgroupplot[title={Direct inverse-speed prediction},xlabel={$r_{d,1000}/3$ / mWb},ylabel={$r_{d,3000}$ / mWb}]
\addplot[measured samples] table[x=predicted_mWb,y=observed_mWb]{figures/data/experimental/common_speed.dat};
\addplot[guide quantity,guide pattern,domain=-0.40551822579603986:0.593869304645285] {x};
\nextgroupplot[title={Equivalent residual at common keys},xlabel={$\Delta R_{\mathrm{eq},1000}$ / m$\Omega$},ylabel={$\Delta R_{\mathrm{eq},3000}$ / m$\Omega$}]
\addplot[measured samples] table[x=req_low_mOhm,y=req_high_mOhm]{figures/data/experimental/common_speed.dat};
\addplot[guide quantity,guide pattern,domain=-7.541175206864074:33.91202647495261] {x};
\end{groupplot}
\end{tikzpicture}
\caption{Two pointwise checks on the same 36 common valid keys at a rotor reference of 70 $^\circ$C. Dashed equality lines are the sole constant-resistance-mismatch predictions. The left plot tests the directly observed residual; the right includes the measured pair-current denominator. The systematic spread away from equality cannot be summarized reliably by a single fitted slope.}
\label{fig:experimental-scaling}
\end{figure}


Measured pair magnitudes across speeds differ by at most
\SI{0.300}{\ampere}. Including their known ratio gives the refined
$H_0$ prediction $r_{d,1000}a_{q,3000}/(3a_{q,1000})$;
its discrepancy RMS remains \SI{0.1571}{\milli\weber}.
On the same keys, the pointwise equivalent difference
$\Delta R_{\mathrm{eq},3000}-\Delta R_{\mathrm{eq},1000}$ has median
$-2.622$ m$\Omega$, raw MAD $1.560$ m$\Omega$ and RMS $4.234$ m$\Omega$,
with range $[-15.178,2.698]$ m$\Omega$. The difference is operating-point
dependent, not just a common displacement between two flat maps.

These are descriptive model checks, not a formal statistical rejection:
there is no calibrated noise model, independent winding-temperature control
verification, or propagated pairing uncertainty. Recorded speed medians
are 999.62 and 2999.78 rpm and the stator reference is 70 $^\circ$C in both
blocks; neither establishes identical magnetic or winding states.
Nevertheless, the spatial and pointwise speed-dependent structure provides
evidence against explaining the observations solely by one constant
resistance mismatch. Additional speed-dependent or operating-point-dependent
effects are required to explain the measurement. Their physical origin is
not established by this comparison.
